% Part: model-theory
% Chapter: models-of-arithmetic
% Section: introduction

\documentclass[../../../include/open-logic-section]{subfiles}

\begin{document}

\olfileid{mod}{mar}{int}
\section{Inleiding}

Het \emph{standaardmodel} van de rekenkunde is de
!!{structure}~$\Struct{N}$ waarvan het domein precies de natuurlijke
getallen is: $\Domain{N} = \Nat$. De tekens $\Obj{0}$, $\prime$, $+$,
$\times$ en $<$ krijgen er hun gewone betekenis. Concreet is $\Obj{0}$ het
getal $0$. Het teken $\prime$ staat voor de functie die van een getal zijn
opvolger maakt. Het teken $+$ betekent optellen en $\times$ betekent
vermenigvuldigen binnen~$\Nat$. In formules:
\begin{align*}
  \Assign{\Obj{0}}{N} & = 0\\
  \Assign{\prime}{N}(n) & = n + 1\\
  \Assign{+}{N}(n, m) & = n + m\\
  \Assign{\times}{N}(n, m) & = nm
\end{align*}.
Maar een structuur voor $\Lang{L_A}$ hoeft niet~$\Nat$ als domein te hebben.
Neem bijvoorbeeld $\Struct{M}$ met
$\Domain{M} = \{a\}^*$. Dit domein bestaat uit alle eindige rijtjes die
alleen het teken~$a$ gebruiken: $\emptyset$, $a$, $aa$, $aaa$, \dots. Geef
de tekens vervolgens deze betekenissen:
\begin{align*}
  \Assign{\Obj{0}}{M} & = \emptyset\\
  \Assign{\prime}{M}(s) & = s \concat a\\
  \Assign{+}{M}(n, m) & = a^{n + m}\\
  \Assign{\times}{M}(n, m) & = a^{nm}
\end{align*}.
De !!{element}s van de !!{domain}s zijn in deze twee structuren anders.
Maar de interpretatiefuncties leggen tussen die elementen precies dezelfde
verbanden. In die zin zijn de structuren ``eigenlijk hetzelfde''. De twee
!!{structure}s zijn, in vaktaal, \emph{isomorf}.

De compactheidsstelling heeft een eenvoudig gevolg. Neem een theorie waarvan
alle zinnen waar zijn in~$\Struct{N}$. Die theorie heeft ook modellen die
niet isomorf zijn met~$\Struct{N}$. Zulke structuren heten
\emph{niet-standaard}. In een standaardmodel---dus in $\Struct{N}$ en in
elke !!{structure} die daarmee isomorf is---bereik je ieder !!{element} met
een standaardnummer~$\num{n}$. Dat is de vaste gesloten term die het
natuurlijke getal n noemt. Daarom geldt
\[
\Domain{N} = \Setabs{\Value{\num{n}}{N}}{n \in \Nat}
\]
In een niet-standaard model lukt dat niet. Als $\Struct{M}$ niet-standaard
is, bevat het minstens één $x \in \Domain{M}$ waarvoor
$x \neq \Value{\num{n}}{M}$ geldt voor ieder~$n$.

Die extra, niet-standaard elementen zijn opvallend: binnen het model gedragen
ze zich als ``oneindige natuurlijke getallen''. Hun bestaan helpt ook te
verklaren waar onvolledigheid vandaan komt. Neem bijvoorbeeld
$\OCon[\Th{PA}]$, de zin die zegt dat je uit de Peano-rekenkunde geen
tegenspraak kunt afleiden. Uitgeschreven is dat $\lnot
\lexists[x][\OPrf[\Th{PA}](x, \gn{\lfalse})]$. De theorie $\Th{PA}$ bewijst
noch $\OCon[\Th{PA}]$, noch $\lnot \OCon[\Th{PA}]$. Je kunt daarom elk van
beide aan $\Th{PA}$ toevoegen zonder een tegenspraak te krijgen. Omdat
$\Th{PA}$ consistent is, geldt $\Sat{N}{\OCon[\Th{PA}]}$, en dus
$\Sat/{N}{\lnot \OCon[\Th{PA}]}$. Het standaardmodel $\Struct{N}$ is dus
\emph{geen} model van $\Th{PA} \cup \{\lnot \OCon[\Th{PA}]\}$. Alle
modellen van die uitgebreidere theorie moeten niet-standaard zijn. Ieder
model van $\Th{PA} \cup \{\lnot \OCon[\Th{PA}]\}$ moet dus een
!!{element} bevatten dat als getuige voor
$\lexists[x][\OPrf[\Th{PA}](\gn{\lfalse})]$ optreedt en die formule in het
model waar maakt. Het element geldt daar als het G\"odelgetal---de
getalcode---van !!a{derivation} van een tegenspraak uit~$\Th{PA}$. Zo'n
!!{element} kan geen standaard element zijn, want $\Th{PA} \Proves \lnot
\OPrf[\Th{PA}](\num{n}, \gn{\lfalse})$ voor elke~$n$.

\end{document}
